\chapter{2009年湖南卷（理）}

\section{单选题}

\begin{problem}
若 \(\log_{2} a < 0\)，\(\left(\frac{1}{2}\right)^b > 1\)，则（\quad）

\choices
    {\(a > 1\)，\(b > 0\)}
    {\(a > 1\)，\(b < 0\)}
    {\(0 < a < 1\)，\(b > 0\)}
    {\(0 < a < 1\)，\(b < 0\)}
\end{problem}

\begin{answer}
D
\end{answer}

\begin{solution}
因为对数的底数 \(2>1\)，所以 \(\log_{2}a<0\) 等价于 \(0<a<1\). 又因为底数 \(\frac{1}{2}\) 在 \((0,1)\) 内，\(\left(\frac{1}{2}\right)^b>1\) 等价于 \(b<0\). 故选 D.
\end{solution}

\begin{problem}
对于非零向量 \(\bs{a}\)，\(\bs{b}\)，“\(\bs{a} + \bs{b} = 0\)”是“\(\bs{a} \parallel \bs{b}\)”的（\quad）

\choices
    {充分不必要条件}
    {必要不充分条件}
    {充分必要条件}
    {既不充分也不必要条件}
\end{problem}

\begin{answer}
A
\end{answer}

\begin{solution}
若 \(\bs{a}+\bs{b}=\bs{0}\)，则 \(\bs{b}=-\bs{a}\)，由于 \(\bs{a}\) 和 \(\bs{b}\) 均非零，它们共线，即 \(\bs{a}\parallel\bs{b}\). 反之，\(\bs{a}\parallel\bs{b}\) 只说明存在实数 \(k\) 使 \(\bs{b}=k\bs{a}\)，不一定有 \(k=-1\)，所以不能推出 \(\bs{a}+\bs{b}=\bs{0}\). 故“\(\bs{a}+\bs{b}=\bs{0}\)”是“\(\bs{a}\parallel\bs{b}\)”的充分不必要条件，选 A.
\end{solution}

\begin{problem}
将函数 \( y = \sin x \) 的图象向左平移 \( \varphi \)（\(0 \leqslant \varphi < 2\pi\)）个单位后，得到函数 \( y = \sin \left(x - \frac{\pi}{6}\right) \) 的图象，则 \( \varphi \) 等于（\quad）

\choices
    {\( \frac{\pi}{6} \)}
    {\( \frac{5\pi}{6} \)}
    {\( \frac{7\pi}{6} \)}
    {\( \frac{11\pi}{6} \)}
\end{problem}

\begin{answer}
D
\end{answer}

\begin{solution}
函数图象向左平移 \(\varphi\) 个单位后，函数变为 \(y=\sin(x+\varphi)\). 要与 \(y=\sin\left(x-\frac{\pi}{6}\right)\) 表示同一函数，必须有 \(\varphi\equiv-\frac{\pi}{6}\pmod{2\pi}\). 结合 \(0\leqslant\varphi<2\pi\)，得 \(\varphi=\frac{11\pi}{6}\)，故选 D.
\end{solution}

\begin{problem}
如图，当参数 \(\lambda = \lambda_{1}\)，\(\lambda_{2}\) 时，连续函数 \(y = \frac{x}{1 + \lambda x}\)（\(x \geqslant 0\)）的图象分别对应曲线 \(C_{1}\) 和 \(C_{2}\)，则（\quad）

\bitmapfigure[max width=0.42\linewidth]{img/2009/hunan_science/q4_fig1.png}

\choices
    {\(0 < \lambda_{1} < \lambda_{2}\)}
    {\(0 < \lambda_{2} < \lambda_{1}\)}
    {\(\lambda_{1} < \lambda_{2} < 0\)}
    {\(\lambda_{2} < \lambda_{1} < 0\)}
\end{problem}

\begin{answer}
B
\end{answer}

\begin{solution}
函数在 \(x\geqslant0\) 上连续，故 \(1+\lambda x\) 不能在该区间内为零，从而 \(\lambda\geqslant0\). 图中两条曲线均不是直线，故 \(\lambda_1\)，\(\lambda_2>0\). 对任意 \(x>0\)，若 \(\lambda_1>\lambda_2\)，则 \(1+\lambda_1x>1+\lambda_2x\)，从而 \(\frac{x}{1+\lambda_1x}<\frac{x}{1+\lambda_2x}\). 图中 \(C_2\) 在 \(C_1\) 的上方，所以 \(\lambda_2<\lambda_1\)，故选 B.
\end{solution}

\begin{problem}
从 \(10\) 名大学毕业生中选 \(3\) 个人担任村长助理，则甲，乙至少有 \(1\) 人入选，而丙没有入选的不同选法的种数为（\quad）

\choices
    {85}
    {56}
    {49}
    {28}
\end{problem}

\begin{answer}
C
\end{answer}

\begin{solution}
丙没有入选时，只能从其余 \(9\) 人中选 \(3\) 人；再去掉甲、乙都没有入选的情形，所求种数为 \(\mathrm C_9^3-\mathrm C_7^3=84-35=49\). 故选 C.
\end{solution}

\begin{problem}
已知 \(D\) 是由不等式组 \(\begin{cases}
x - 2y \geqslant 0\\
x + 3y \geqslant 0
\end{cases}\) 所确定的平面区域. 则圆 \(x^2 + y^2 = 4\) 在区域 \(D\) 内的弧长为（\quad）

\choices
    {\( \frac{\pi}{4} \)}
    {\( \frac{\pi}{2} \)}
    {\( \frac{3\pi}{4} \)}
    {\( \frac{3\pi}{2} \)}
\end{problem}

\begin{answer}
B
\end{answer}

\begin{solution}
区域 \(D\) 满足 \(y\leqslant\frac{x}{2}\) 且 \(y\geqslant-\frac{x}{3}\)，从而 \(-\frac{x}{3}\leqslant\frac{x}{2}\)，故 \(x\geqslant0\). 因此，\(D\) 是由两条射线 \(y=\frac{x}{2}\) 与 \(y=-\frac{x}{3}\) 围成的角域. 设其圆心角为 \(\theta\)，则 \(\tan\theta=\frac{\frac12+\frac13}{1-\frac12\cdot\frac13}=1\). 两条射线夹角为锐角，故 \(\theta=\frac{\pi}{4}\). 圆半径为 \(2\)，弧长为 \(2\theta=\frac{\pi}{2}\)，故选 B.
\end{solution}

\begin{problem}
正方体 \(ABCD - A_{1}B_{1}C_{1}D_{1}\) 的棱上到异面直线 \(AB\)，\(CC_{1}\) 的距离相等的点的个数为（\quad）

\choices
    {2}
    {3}
    {4}
    {5}
\end{problem}

\begin{answer}
C
\end{answer}

\begin{solution}
设正方体棱长为 \(1\)，取坐标 \(A(0,0,0)\)，\(B(1,0,0)\)，\(C(1,1,0)\)，\(D(0,1,0)\)，\(A_{1}(0,0,1)\)，则直线 \(AB\) 为 \(y=z=0\)，直线 \(CC_{1}\) 为 \(x=y=1\). 对任意点 \(P(x,y,z)\)，到两直线的距离平方分别为 \(y^{2}+z^{2}\) 和 \((x-1)^{2}+(y-1)^{2}\)，所以所求条件为
\[
y^2+z^2=(x-1)^2+(y-1)^2
\]
逐条考察各棱，其中参数 \(t\) 均取遍 \([0,1]\). 在 \(AB\)、\(BC\)、\(CD\)、\(DA\) 上分别取点 \((t,0,0)\)、\((1,t,0)\)、\((t,1,0)\)、\((0,t,0)\)，代入上式分别得到 \(0=(t-1)^2+1\)、\(t^2=(t-1)^2\)、\(1=(t-1)^2\)、\(t^2=1+(t-1)^2\). 因而 \(AB\) 上无解，\(BC\) 上 \(t=\frac12\)，即其中点，\(CD\) 上 \(t=0\)，即 \(D\)，\(DA\) 上 \(t=1\)，即 \(D\).

在 \(AA_{1}\)、\(BB_{1}\)、\(CC_{1}\)、\(DD_{1}\) 上分别取点 \((0,0,t)\)、\((1,0,t)\)、\((1,1,t)\)、\((0,1,t)\)，代入上式分别得到 \(t^2=2\)、\(t^2=1\)、\(1+t^2=t^2\)、\(1+t^2=1\). 因而 \(AA_{1}\)、\(CC_{1}\) 上无解，\(BB_{1}\) 上 \(t=1\)，即 \(B_{1}\)，\(DD_{1}\) 上 \(t=0\)，即 \(D\).

在 \(A_{1}B_{1}\)、\(B_{1}C_{1}\)、\(C_{1}D_{1}\)、\(D_{1}A_{1}\) 上分别取点 \((t,0,1)\)、\((1,t,1)\)、\((t,1,1)\)、\((0,t,1)\)，代入上式分别得到 \(1=(t-1)^2+1\)、\(t^2+1=(t-1)^2\)、\(2=(t-1)^2\)、\(t^2+1=1+(t-1)^2\). 因而 \(A_{1}B_{1}\) 上 \(t=1\)，即 \(B_{1}\)，\(B_{1}C_{1}\) 上 \(t=0\)，即 \(B_{1}\)，\(C_{1}D_{1}\) 上无解，\(D_{1}A_{1}\) 上 \(t=\frac12\)，即其中点.

去掉重复端点后，满足条件的点为 \(D\)、\(B_{1}\)、\(BC\) 的中点和 \(D_{1}A_{1}\) 的中点，共 \(4\) 个，故选 C.
\end{solution}

\begin{problem}
设函数 \( y = f(x) \) 在 \((- \infty, + \infty)\) 内有定义. 对于给定的正数 \(K\)，定义函数 \( f_{K}(x) = \begin{cases}
f(x), & f(x) \leqslant K,\\
K, & f(x) > K.
\end{cases} \) 取函数 \( f(x) = 2 - x - \e^{-x} \). 若对任意的 \( x \in (-\infty, +\infty)\)，恒有 \( f_{K}(x) = f(x) \)，则（\quad）

\choices
    {\(K\) 的最大值为2}
    {\(K\) 的最小值为2}
    {\(K\) 的最大值为1}
    {\(K\) 的最小值为1}
\end{problem}

\begin{answer}
D
\end{answer}

\begin{solution}
对 \(f(x)=2-x-\e^{-x}\) 求导，得 \(f'(x)=\e^{-x}-1\). 当 \(x<0\) 时 \(f'(x)>0\)，当 \(x>0\) 时 \(f'(x)<0\)，所以 \(f(x)\) 在 \(x=0\) 处取得最大值 \(f(0)=1\). 函数 \(f_K(x)\) 恒等于 \(f(x)\) 的充要条件是 \(K\geqslant1\)，故 \(K\) 的最小值为 \(1\)，选 D.
\end{solution}

\section{填空题}

\begin{problem}
某班共 \(30\) 人，其中 \(15\) 人喜爱篮球运动，\(10\) 人喜爱乒乓球运动，\(8\) 人对这两项运动都不喜爱，则喜爱篮球运动但不喜爱乒乓球运动的人数为\fillinblank{}.
\end{problem}

\begin{answer}
\(12\)
\end{answer}

\begin{solution}
喜爱至少一项运动的人数为 \(30-8=22\). 设同时喜爱两项运动的人数为 \(t\)，则 \(15+10-t=22\)，解得 \(t=3\). 所以只喜爱篮球运动的人数为 \(15-3=12\).
\end{solution}

\begin{problem}
在 \((1 + x)^{3} + (1 + \sqrt{x})^{3} + (1 + \sqrt[3]{x})^{3}\) 的展开式中，\(x\) 的系数为\fillinblank{}.
\end{problem}

\begin{answer}
\(7\)
\end{answer}

\begin{solution}
\((1+x)^3\) 中 \(x\) 的系数为 \(3\)，\((1+\sqrt{x})^3=1+3\sqrt{x}+3x+x\sqrt{x}\) 中 \(x\) 的系数为 \(3\)，\((1+\sqrt[3]{x})^3=1+3\sqrt[3]{x}+3\sqrt[3]{x^2}+x\) 中 \(x\) 的系数为 \(1\). 三者相加，\(x\) 的系数为 \(3+3+1=7\).
\end{solution}

\begin{problem}
若 \(x \in \left(0, \frac{\pi}{2}\right)\)，则 \(2\tan x + \tan \left(\frac{\pi}{2} - x\right)\) 的最小值为\fillinblank{}.
\end{problem}

\begin{answer}
\(2\sqrt{2}\)
\end{answer}

\begin{solution}
因为 \(x\in\left(0,\frac{\pi}{2}\right)\)，令 \(t=\tan x\)，则 \(t>0\)，且 \(\tan\left(\frac{\pi}{2}-x\right)=\frac1t\). 因此原式为 \(2t+\frac1t\). 由基本不等式，\(2t+\frac1t\geqslant2\sqrt{2t\cdot\frac1t}=2\sqrt2\)，等号在 \(2t=\frac1t\)，即 \(t=\frac1{\sqrt2}\) 时成立，故最小值为 \(2\sqrt2\).
\end{solution}

\begin{problem}
已知以双曲线 \(C\) 的两个焦点及虚轴的两个端点为顶点的四边形中，有一个内角为 \( 60^{\circ} \)，则双曲线 \(C\) 的离心率为\fillinblank{}.
\end{problem}

\begin{answer}
\(\frac{\sqrt{6}}{2}\)
\end{answer}

\begin{solution}
设双曲线方程为 \(\frac{x^2}{a^2}-\frac{y^2}{b^2}=1\)，焦点为 \((\pm c,0)\)，虚轴端点为 \((0,\pm b)\)，其中 \(c^2=a^2+b^2\). 在焦点处的内角 \(\theta\) 满足 \(\cos\theta=\frac{c^2-b^2}{c^2+b^2}\). 由于双曲线中 \(c>b\)，题设的 \(60^\circ\) 角只能是焦点处的锐角，故 \(\frac{c^2-b^2}{c^2+b^2}=\frac12\)，从而 \(c^2=3b^2\). 于是 \(a^2=c^2-b^2=2b^2\)，所以离心率 \(e=\frac ca=\sqrt{\frac32}=\frac{\sqrt6}{2}\).
\end{solution}

\begin{problem}
一个总体分为 \(A\)，\(B\) 两层，其个体数之比为 \(4:1\)，用分层抽样方法从总体中抽取一个容量为 \(10\) 的样本. 已知 \(B\) 层中甲，乙都被抽到的概率为 \(\frac{1}{28}\)，则总体中的个体数为\fillinblank{}.
\end{problem}

\begin{answer}
\(40\)
\end{answer}

\begin{solution}
按 \(4:1\) 分层抽样时，容量为 \(10\) 的样本中 \(B\) 层应抽取 \(2\) 人. 设 \(B\) 层共有 \(N_B\) 人，则甲、乙同时被抽到的概率为 \(\frac{1}{\mathrm C_{N_B}^{2}}=\frac1{28}\)，所以 \(N_B(N_B-1)=56\)，解得 \(N_B=8\). 总体人数为 \(5N_B=40\).
\end{solution}

\begin{problem}
在半径为 \(13\) 的球面上有 \(A\)，\(B\)，\(C\) 三点，\(AB = 6\)，\(BC = 8\)，\(CA = 10\)，则
\begin{enumerate}
\item 球心到平面 \(ABC\) 的距离为\fillinblank{}；
\item 过 \(A\)，\(B\) 两点的大圆面与平面 \(ABC\) 所成二面角为 （锐角） 的正切值为\fillinblank{}.
\end{enumerate}
\end{problem}

\begin{answer}
\(12\)，\(3\)
\end{answer}

\begin{solution}
\begin{enumerate}
\item 因为 \(AB^2+BC^2=6^2+8^2=10^2=AC^2\)，所以 \(\triangle ABC\) 在 \(B\) 处为直角三角形. 其外接圆半径为斜边 \(AC\) 的一半，即 \(5\). 设球心到平面 \(ABC\) 的距离为 \(h\)，则 \(13^2=h^2+5^2\)，解得 \(h=12\).
\item 建立直角坐标系，令 \(B=(0,0,0)\)，\(A=(6,0,0)\)，\(C=(0,8,0)\). 平面 \(ABC\) 为 \(z=0\)，球心可取 \(O=(3,4,12)\). 平面 \(OAB\) 是过 \(A\)，\(B\) 两点的大圆面，其一个法向量为 \((0,-3,1)\)；平面 \(ABC\) 的法向量为 \((0,0,1)\). 设两平面锐角为 \(\theta\)，则 \(\cos\theta=\frac{1}{\sqrt{10}}\)，所以 \(\tan\theta=3\).
\end{enumerate}
\end{solution}

\begin{problem}
将正 \(\triangle ABC\) 分割成 \(n^2\)（\(n \geqslant 2\)，\(n \in \N^*\)）个全等的小正三角形（图1，图2分别给出了 \(n = 2\)，\(3\) 的情形）. 在每个三角形的顶点各放置一个数，使位于 \(\triangle ABC\) 的三边及平行于某边的任一直线上的数（当数的个数不少于\(3\)时）都分别依次成等差数列. 若顶点 \(A\)，\(B\)，\(C\) 处的三个数互不相同且和为 \(1\)，记所有顶点上的数的和为 \(f(n)\)，则有 \(f(2) = 2\)，\(f(3) =\fillinblank{}\)，\(\cdots\)，\(f(n) =\fillinblank{}\).

\bitmapfigure[max width=0.82\linewidth]{img/2009/hunan_science/q15_fig12.png}
\end{problem}

\begin{answer}
\(\frac{10}{3}\)，\(\frac{1}{6}(n + 1)(n + 2)\)
\end{answer}

\begin{solution}
设 \(A\)，\(B\)，\(C\) 三个顶点上的数分别为 \(a\)，\(b\)，\(c\)，则 \(a+b+c=1\). 把网格顶点记为 \(P_{ijk}\)，其中 \(i\)，\(j\)，\(k\geqslant0\)，\(i+j+k=n\)，并约定 \(i\)，\(j\)，\(k\) 分别对应顶点 \(A\)，\(B\)，\(C\) 的重数. 沿边 \(AB\) 与 \(AC\) 的等差条件先给出边上点的数值：对应于 \((i,j,0)\) 的点数值为 \(\frac{ia+jb}{n}\)，对应于 \((i,0,k)\) 的点数值为 \(\frac{ia+kc}{n}\). 对固定的 \(i\)，所有满足 \(i+j+k=n\) 的点在一条平行于 \(BC\) 的直线上，该直线两端的数值已知，且这些数依次成等差数列，所以中间点的数值必为 \(\frac{ia+jb+kc}{n}\). 因此全部网格顶点的数值均可写成 \(\frac{ia+jb+kc}{n}\). 网格顶点总数为 \(N=\frac{(n+1)(n+2)}2\). 由指标 \(i\)，\(j\)，\(k\) 的对称性，\(\sum i=\sum j=\sum k=\frac{nN}{3}\)，故
\[
f(n)=\sum_{i+j+k=n}\frac{ia+jb+kc}{n}
=\frac{(a+b+c)N}{3}
=\frac{(n+1)(n+2)}6
\]
代入 \(n=3\)，得 \(f(3)=\frac{10}{3}\)，且 \(f(2)=2\) 与题面一致.
\end{solution}

\section{解答题}

\begin{problem}
在 \(\triangle ABC\) 中，已知 \(2\overrightarrow{AB} \cdot \overrightarrow{AC} = \sqrt{3} |\overrightarrow{AB}| \cdot |\overrightarrow{AC}| = 3BC^2\)，求角 \(A\)，\(B\)，\(C\) 的大小.
\end{problem}

\begin{solution}
由 \(2\overrightarrow{AB}\cdot\overrightarrow{AC}=\sqrt3|\overrightarrow{AB}||\overrightarrow{AC}|\)，得 \(\cos A=\frac{\overrightarrow{AB}\cdot\overrightarrow{AC}}{|\overrightarrow{AB}||\overrightarrow{AC}|}=\frac{\sqrt3}{2}\)，所以 \(A=\frac{\pi}{6}\). 记 \(AB=u\)，\(AC=v\)，\(BC=w\)，由题设 \(3w^2=\sqrt3uv\)，即 \(w^2=\frac{uv}{\sqrt3}\). 另一方面由余弦定理，\(w^2=u^2+v^2-\sqrt3uv\)，故 \(u^2+v^2=\frac4{\sqrt3}uv\). 令 \(r=\frac uv>0\)，则 \(r+\frac1r=\frac4{\sqrt3}\)，即 \(3r^2-4\sqrt3r+3=0\)，解得 \(r=\sqrt3\) 或 \(r=\frac1{\sqrt3}\). 当 \(u=\sqrt3v\) 时，\(w^2=v^2\)，三边之比为 \(\sqrt3:1:1\)，所以 \(A=B=\frac{\pi}{6}\)，由三角形内角和得 \(C=\frac{2\pi}{3}\)，即 \((A,B,C)=\left(\frac{\pi}{6},\frac{\pi}{6},\frac{2\pi}{3}\right)\). 当 \(u=\frac{v}{\sqrt3}\) 时，\(w^2=u^2\)，三边之比为 \(1:\sqrt3:1\)，所以 \(A=C=\frac{\pi}{6}\)，\(B=\frac{2\pi}{3}\)，即 \((A,B,C)=\left(\frac{\pi}{6},\frac{2\pi}{3},\frac{\pi}{6}\right)\).
\end{solution}

\begin{problem}
为拉动经济增长，某市决定新建一批重点工程，分别为基础设施工程，民生工程和产业建设工程三类. 这三类工程所含项目的个数分别占总数的 \(\frac{1}{2}\)，\(\frac{1}{3}\)，\(\frac{1}{6}\). 现在 \(3\) 名工人独立地从中任选一个项目参与建设.
\begin{enumerate}
\item 求他们选择的项目所属类别互不相同的概率；
\item 记 \( \xi \) 为 \(3\) 人中选择的项目属于基础设施工程或产业建设工程的人数，求 \( \xi \) 的分布列及数学期望.
\end{enumerate}
\end{problem}

\begin{solution}
\begin{enumerate}
\item 三名工人所选类别互不相同的排列有 \(3!\) 种，每一种指定排列的概率为 \(\frac12\cdot\frac13\cdot\frac16\)，所以所求概率为 \(3!\cdot\frac12\cdot\frac13\cdot\frac16=\frac16\).
\item 每名工人选择基础设施工程或产业建设工程的概率为 \(p=\frac12+\frac16=\frac23\)，选择民生工程的概率为 \(1-p=\frac13\). 三名工人独立选择，因此 \(\xi\) 服从二项分布，且 \(P(\xi=k)=\mathrm C_3^k\left(\frac23\right)^k\left(\frac13\right)^{3-k}\)，\(k=0\)，\(1\)，\(2\)，\(3\). 具体分布列为
\begin{center}
\begin{tabular}{c|cccc}
\(\xi\) & \(0\) & \(1\) & \(2\) & \(3\)\\
\hline
\(P\) & \(\frac1{27}\) & \(\frac2{9}\) & \(\frac4{9}\) & \(\frac8{27}\)
\end{tabular}
\end{center}
数学期望为 \(E(\xi)=3p=3\cdot\frac23=2\).
\end{enumerate}
\end{solution}

\begin{problem}
如图，在正三棱柱 \(ABC - A_{1}B_{1}C_{1}\) 中，\(AB = \sqrt{2}AA_{1}\)，点 \(D\) 是 \(A_{1}B_{1}\) 的中点，点 \(E\) 在 \(A_{1}C_{1}\) 上，且 \(DE \perp AE\).
\begin{enumerate}
\item 证明：平面 \(ADE\) \(\perp\) 平面 \(ACC_{1}A_{1}\)；
\item 求直线 \(AD\) 和平面 \(ABC_{1}\) 所成角的正弦值.
\end{enumerate}

\bitmapfigure[max width=0.45\linewidth]{img/2009/hunan_science/q18_fig1.png}
\end{problem}

\begin{solution}
\begin{enumerate}
\item 正三棱柱的侧棱 \(AA_1\) 垂直于上底面 \(A_1B_1C_1\)，而 \(DE\) 在上底面内，所以 \(DE\perp AA_1\). 又已知 \(DE\perp AE\). 由题图中 \(E\) 为 \(A_1C_1\) 的内点，\(AE\) 与 \(AA_1\) 是平面 \(ACC_1A_1\) 内相交的两条不同直线，因此 \(DE\perp\) 平面 \(ACC_1A_1\). 由于 \(DE\) 在平面 \(ADE\) 内，故平面 \(ADE\perp\) 平面 \(ACC_1A_1\).
\item 令 \(AA_1=1\)，则 \(AB=\sqrt2\). 建立空间直角坐标系，取 \(A=(0,0,0)\)，\(B=(\sqrt2,0,0)\)，\(C=\left(\frac{\sqrt2}{2},\frac{\sqrt6}{2},0\right)\)，\(A_1=(0,0,1)\)，\(C_1=\left(\frac{\sqrt2}{2},\frac{\sqrt6}{2},1\right)\)，\(D=\left(\frac{\sqrt2}{2},0,1\right)\). 于是直线 \(AD\) 的方向向量为 \(\left(\frac{\sqrt2}{2},0,1\right)\)，平面 \(ABC_1\) 的一个法向量为 \(\overrightarrow{AB}\times\overrightarrow{AC_1}=(0,-\sqrt2,\sqrt3)\). 设直线 \(AD\) 与平面 \(ABC_1\) 所成的角为 \(\theta\)，则
\[
\sin\theta=\frac{\left|\left(\frac{\sqrt2}{2},0,1\right)\cdot(0,-\sqrt2,\sqrt3)\right|}{\sqrt{\frac12+1}\sqrt{2+3}}=\frac{\sqrt3}{\frac{\sqrt6}{2}\sqrt5}=\frac{\sqrt{10}}5
\]
\end{enumerate}
\end{solution}

\begin{problem}
某地建一座桥，两端的桥墩已建好，这两墩相距 \(m\text{ 米}\)，余下工程只需要建两端桥墩之间的桥面和桥墩. 经测算，一个桥墩的工程费用为 \(256\text{ 万元}\)，距离为 \(x\text{ 米}\) 的相邻两墩之间的桥面工程费用为 \((2 + \sqrt{x})x\text{ 万元}\). 假设桥墩等距离分布，所有桥墩都视为点，且不考虑其他因素，记余下工程的费用为 \(y\text{ 万元}\).
\begin{enumerate}
\item 试写出 \(y\) 关于 \(x\) 的函数关系式；
\item 当 \(m = 640\text{ 米}\) 时，需新建多少个桥墩才能使 \(y\) 最小？
\end{enumerate}
\end{problem}

\begin{solution}
\begin{enumerate}
\item 设新建桥墩数为 \(N\)，则相邻桥墩之间的桥面共有 \(N+1\) 段，所以 \((N+1)x=m\)，即 \(N=\frac mx-1\). 因此
\[
y=256\left(\frac mx-1\right)+\frac mx(2+\sqrt{x})x=\frac{256m}{x}+m\sqrt{x}+2m-256
\]
其中 \(0<x\leqslant m\)，实际等距分布时还需满足 \(\frac mx\in\N^*\).
\item 当 \(m=640\) 时，令 \(y(x)=\frac{163840}{x}+640\sqrt{x}+1024\). 在 \(0<x\leqslant640\) 上，\(y'(x)=-\frac{163840}{x^2}+\frac{320}{\sqrt{x}}\). 令 \(y'(x)=0\)，得 \(x^{3/2}=512\)，所以 \(x=64\). 当 \(0<x<64\) 时 \(y'(x)<0\)，当 \(64<x\leqslant640\) 时 \(y'(x)>0\)，故 \(y\) 在 \(x=64\) 处取得最小值. 此时桥面段数为 \(\frac{640}{64}=10\)，需新建桥墩数为 \(10-1=9\).
\end{enumerate}
\end{solution}

\begin{problem}
在平面直角坐标系 \(xOy\) 中，点 \(P\) 到点 \(F(3,0)\) 的距离的 \(4\) 倍与它到直线 \(x = 2\) 的距离的 \(3\) 倍之和记为 \(d\). 当点 \(P\) 运动时，\(d\) 恒等于点 \(P\) 的横坐标与 \(18\) 之和.
\begin{enumerate}
\item 求点 \(P\) 的轨迹 \(C\)；
\item 设过点 \(F\) 的直线 \(l\) 与轨迹 \(C\) 相交于 \(M\)，\(N\) 两点，求线段 \(MN\) 长度的最大值.
\end{enumerate}
\end{problem}

\begin{solution}
\begin{enumerate}
\item 设 \(P=(x,y)\)，则 \(d=4\sqrt{(x-3)^2+y^2}+3|x-2|\). 当 \(x\geqslant2\) 时，原条件化为 \(4\sqrt{(x-3)^2+y^2}=24-2x\)，平方并整理得 \(3x^2+4y^2=108\)，即 \(\frac{x^2}{36}+\frac{y^2}{27}=1\)，并取其 \(x\geqslant2\) 的部分. 当 \(x\leqslant2\) 时，原条件化为 \(4\sqrt{(x-3)^2+y^2}=4x+12\)，平方并整理得 \(y^2=12x\)，并由 \(y^2\geqslant0\) 得 \(0\leqslant x\leqslant2\). 故
\[
C:\begin{cases}
\dfrac{x^2}{36}+\dfrac{y^2}{27}=1, & x\geqslant2,\\
y^2=12x, & 0\leqslant x\leqslant2.
\end{cases}
\]
\item 先考虑水平直线 \(y=0\). 它与抛物线弧交于 \((0,0)\)，与椭圆弧交于 \((6,0)\)，故此时 \(MN=6\). 对于非水平且非竖直直线 \(l\)，设其方程为 \(y=k(x-3)\)，记 \(s=\sqrt{1+k^2}>1\)，再令 \(q=k^2=s^2-1\). 将直线代入椭圆方程，得到 \((3+4q)x^2-24qx+36q-108=0\)，其两个根之差为 \(\frac{36s}{3+4q}\). 将直线代入抛物线方程，得到 \(qx^2-(6q+12)x+9q=0\)，其较小根为 \(x_p=3-\frac6{s+1}\).

当 \(1<s<5\) 时，\(0<x_p<2\)，抛物线给出一个轨迹点；椭圆的另一个交点的横坐标小于 \(2\)，所以椭圆只给出较大的根
\[
x_e=\frac{12q+18s}{3+4q}=\frac{6(s+2)}{2s+1}
\]
因此
\[
MN=s(x_e-x_p)=\frac{3s(7s+5)}{(2s+1)(s+1)}=\frac{3(7s^2+5s)}{2s^2+3s+1}
\]
其导数为 \(\frac{3(11s^2+14s+5)}{(2s^2+3s+1)^2}>0\)，故该表达式在 \(1<s<5\) 上递增，并在 \(s=5\) 的接合位置连续取到 \(\frac{100}{11}\).

当 \(s\geqslant5\) 时，椭圆的两个交点均满足 \(x\geqslant2\). 当 \(s=5\) 时，抛物线弧的交点为 \(x=2\)，恰与椭圆弧的端点重合；当 \(s>5\) 时，抛物线弧没有交点. 因而
\[
MN=\frac{36s^2}{4s^2-1}
\]
这是 \(s\geqslant5\) 上的减函数，在 \(s=5\) 时也等于 \(\frac{100}{11}\). 竖直直线 \(x=3\) 与椭圆的弦长为 \(9<\frac{100}{11}\)，所以线段 \(MN\) 长度的最大值为 \(\frac{100}{11}\).
\end{enumerate}
\end{solution}

\begin{problem}
对于数列 \(\{u_n\}\)，若存在常数 \(M > 0\)，对任意的 \(n \in \N^*\)，恒有 \(|u_{n+1} - u_n| + |u_n - u_{n-1}| + \cdots + |u_2 - u_1| \leqslant M\)，则称数列 \(\{u_n\}\) 为 \(B\)-数列.
\begin{enumerate}
\item 首项为 \(1\)，公比为 \(q\)（\(|q| < 1\)）的等比数列是否为 \(B\)-数列？请说明理由；
\item 设 \(S_{n}\) 是数列 \(\{x_{n}\}\) 的前 \(n\) 项和. 给出下列两组论断：
\begin{circlelist}
\item A 组：数列 \(\{x_{n}\}\) 是 \(B\)-数列；
\item A 组：数列 \(\{x_{n}\}\) 不是 \(B\)-数列；
\item B 组：数列 \(\{S_{n}\}\) 是 \(B\)-数列；
\item B 组：数列 \(\{S_{n}\}\) 不是 \(B\)-数列.
\end{circlelist}
请以其中一组中的一个论断为条件，另一组中的一个论断为结论组成一个命题. 判断所给命题的真假，并证明你的结论；
\item 若数列 \(\{a_{n}\}\)，\(\{b_{n}\}\) 都是 \(B\)-数列，证明：数列 \(\{a_{n}b_{n}\}\) 也是 \(B\)-数列.
\end{enumerate}
\end{problem}

\begin{solution}
\begin{enumerate}
\item 设 \(u_n=q^{n-1}\). 则
\[
\sum_{k=1}^{n}|u_{k+1}-u_k|=|1-q|\sum_{k=0}^{n-1}|q|^k=|1-q|\frac{1-|q|^n}{1-|q|}<\frac{|1-q|}{1-|q|}
\]
右端与 \(n\) 无关，所以取 \(M=\frac{|1-q|}{1-|q|}\) 即可，故该等比数列是 \(B\)-数列.
\item 取命题“若数列 \(\{x_n\}\) 是 \(B\)-数列，则数列 \(\{S_n\}\) 是 \(B\)-数列”. 令 \(x_n=1\)，则 \(\{x_n\}\) 的各项相等，因而是 \(B\)-数列；但 \(S_n=n\)，从而 \(\sum_{k=1}^{n}|S_{k+1}-S_k|=\sum_{k=1}^{n}1=n\)，不可能被某个常数 \(M\) 对所有 \(n\) 统一界住，所以该命题是假命题.

再取命题“若数列 \(\{S_n\}\) 是 \(B\)-数列，则数列 \(\{x_n\}\) 是 \(B\)-数列”. 若 \(\{S_n\}\) 是 \(B\)-数列，则存在 \(M>0\)，使 \(\sum_{k=1}^{n}|S_{k+1}-S_k|\leqslant M\). 由于 \(S_{k+1}-S_k=x_{k+1}\)，有 \(\sum_{k=1}^{n}|x_{k+1}|\leqslant M\). 于是
\[
\sum_{k=1}^{n}|x_{k+1}-x_k|\leqslant |x_1|+|x_{n+1}|+2\sum_{k=2}^{n}|x_k|\leqslant |x_1|+2M
\]
右端与 \(n\) 无关，所以 \(\{x_n\}\) 是 \(B\)-数列，该命题为真.
\item 由 \(\{a_n\}\) 是 \(B\)-数列，存在 \(M_1>0\)，使其总变差不超过 \(M_1\)，故对任意 \(n\)，\(|a_n|\leqslant |a_1|+M_1=:K_1\). 由 \(\{b_n\}\) 是 \(B\)-数列，存在 \(M_2>0\)，使其总变差不超过 \(M_2\)，故对任意 \(n\)，\(|b_n|\leqslant |b_1|+M_2=:K_2\). 对每个 \(k\)，
\[
|a_{k+1}b_{k+1}-a_kb_k|\leqslant |b_{k+1}||a_{k+1}-a_k|+|a_k||b_{k+1}-b_k|\leqslant K_2|a_{k+1}-a_k|+K_1|b_{k+1}-b_k|
\]
对 \(k=1\)，\(\ldots\)，\(n\) 求和，得
\[
\sum_{k=1}^{n}|a_{k+1}b_{k+1}-a_kb_k|\leqslant K_2M_1+K_1M_2
\]
右端与 \(n\) 无关，故 \(\{a_nb_n\}\) 是 \(B\)-数列.
\end{enumerate}
\end{solution}
